\documentclass[12pt,a4paper]{article}

\usepackage[a4paper,margin=1in]{geometry}
\usepackage{fontspec}
\usepackage{ctex}
\usepackage{amsmath,amssymb}
\usepackage{graphicx}
\usepackage{booktabs}
\usepackage{array}
\usepackage{longtable}
\usepackage{caption}
\usepackage{float}
\usepackage{setspace}
\usepackage{fancyhdr}
\usepackage{hyperref}
\usepackage{xcolor}
\usepackage{enumitem}
\usepackage{listings}
\usepackage{calligra}

\setcounter{tocdepth}{2}
\setmainfont{Noto Serif}
\setmonofont{DejaVu Sans Mono}
\onehalfspacing
\setlength{\parindent}{2em}
\setlength{\parskip}{0.4em}
\pagestyle{fancy}
\fancyhf{}
\fancyfoot[C]{\thepage}
\renewcommand{\headrulewidth}{0pt}
\hypersetup{colorlinks=true,linkcolor=blue,urlcolor=blue,citecolor=blue}
\lstset{basicstyle=\ttfamily\small,breaklines=true,frame=single,columns=fullflexible,showstringspaces=false}
\renewcommand{\contentsname}{Contents}

\newcommand{\HomeworkTitle}{Computational Physics Homework 7 Report}
\newcommand{\StudentID}{2024141230044}
\newcommand{\StudentName}{Fan Yi}
\newcommand{\PhotoFile}{photo.jpg}
\newcommand{\PhotoWidth}{35mm}
\newcommand{\PhotoHeight}{49mm}

\begin{document}

\noindent
\begin{minipage}[c]{0.72\textwidth}
    \begin{center}
        {\LARGE \textbf{\HomeworkTitle}}\\[1.2em]
        {\large Student ID: \StudentID}\\[0.5em]
        {\large Name: \StudentName}
    \end{center}
\end{minipage}%
\hfill
\begin{minipage}[c]{\PhotoWidth}
    \centering
    \IfFileExists{\PhotoFile}{%
        \includegraphics[width=\PhotoWidth,height=\PhotoHeight,keepaspectratio]{\PhotoFile}%
    }{%
        \fbox{\parbox[c][\PhotoHeight][c]{\PhotoWidth}{\centering 2-inch\\ID photo}}%
    }
\end{minipage}
\vspace{1em}
\noindent\rule{\textwidth}{0.8pt}
\vspace{1.5em}

This report summarizes the four homework tasks shown on the last pages of \textit{Computational Physics Lecture 7}. The folder contains several drafts and executable files, so I first identified the versions that are most likely to be the final submissions, and then organized their mathematical results and numerical outputs into one English report.

\newpage
\tableofcontents
\newpage

\section{Selected Source Files}

The working directory contains several versions of the homework. The files chosen as the most likely final drafts are listed in Table~\ref{tab:selected}. The main criterion is whether a file matches the lecture task directly and contains the complete required workflow instead of only a partial experiment.

\begin{table}[H]
\centering
\caption{Selected final drafts used in this report.}
\label{tab:selected}
\begin{tabular}{>{\raggedright\arraybackslash}p{0.13\textwidth} >{\raggedright\arraybackslash}p{0.27\textwidth} >{\raggedright\arraybackslash}p{0.48\textwidth}}
\toprule
Task & Selected file(s) & Reason \\
\midrule
Problem 1 & \texttt{hw7\_1/main.tex}, \texttt{hw7\_1/main.pdf} & Clean electronic version of the handwritten derivation. The file \texttt{手写稿.pdf} is treated as the original handwritten draft rather than the final submission. \\
Problem 2 & \texttt{hw7\_2/main.py} & Covers the full assignment: table-data approximation, error sweep for many $h$ values, and figure generation. The C file only covers the second part. \\
Problem 3 & \texttt{hw7\_3/main.py} & Only complete implementation in the folder, including both the table-data test and the $h$-dependence study. \\
Problem 4(a) & \texttt{hw7\_4/a/a.c} & Direct implementation of the first derivative program requested in the lecture. \\
Problem 4(b) & \texttt{hw7\_4/b/b.c} & Direct implementation of the multi-$\delta$ accuracy test. \\
Problem 4(c) & \texttt{hw7\_4/c/c.py} & Matches the lecture prompt most directly by reproducing the forward-difference error plot for $f(x)=\sin x$. The files \texttt{c+.py} and \texttt{c++.py} appear to be later extensions to higher-order formulas. \\
\bottomrule
\end{tabular}
\end{table}

\section{Theoretical Analysis}

\subsection{Problem 1: finite-difference formulas}

Starting from the shift operator $E f(x)=f(x+h)$, the forward, backward, and central differences are
\[
\Delta = E-1, \qquad \nabla = 1-E^{-1}, \qquad \delta = E^{1/2}-E^{-1/2}.
\]
Expanding these operators and combining them with Taylor series gives the derivative formulas below.

\subsubsection{Forward differences}

\[
f''(x) \approx \frac{f(x+2h)-2f(x+h)+f(x)}{h^2},
\qquad O(h).
\]
\[
f^{(3)}(x) \approx \frac{f(x+3h)-3f(x+2h)+3f(x+h)-f(x)}{h^3},
\qquad O(h).
\]
\[
f^{(4)}(x) \approx \frac{f(x+4h)-4f(x+3h)+6f(x+2h)-4f(x+h)+f(x)}{h^4},
\qquad O(h).
\]
\[
f^{(5)}(x) \approx \frac{f(x+5h)-5f(x+4h)+10f(x+3h)-10f(x+2h)+5f(x+h)-f(x)}{h^5},
\qquad O(h).
\]

\subsubsection{Backward differences}

\[
f''(x) \approx \frac{f(x)-2f(x-h)+f(x-2h)}{h^2},
\qquad O(h).
\]
\[
f^{(3)}(x) \approx \frac{f(x)-3f(x-h)+3f(x-2h)-f(x-3h)}{h^3},
\qquad O(h).
\]
\[
f^{(4)}(x) \approx \frac{f(x)-4f(x-h)+6f(x-2h)-4f(x-3h)+f(x-4h)}{h^4},
\qquad O(h).
\]
\[
f^{(5)}(x) \approx \frac{f(x)-5f(x-h)+10f(x-2h)-10f(x-3h)+5f(x-4h)-f(x-5h)}{h^5},
\qquad O(h).
\]

\subsubsection{Central differences}

\[
f''(x) \approx \frac{f(x+h)-2f(x)+f(x-h)}{h^2},
\qquad O(h^2).
\]
\[
f^{(3)}(x) \approx \frac{f(x+2h)-2f(x+h)+2f(x-h)-f(x-2h)}{2h^3},
\qquad O(h^2).
\]
\[
f^{(4)}(x) \approx \frac{f(x+2h)-4f(x+h)+6f(x)-4f(x-h)+f(x-2h)}{h^4},
\qquad O(h^2).
\]
\[
f^{(5)}(x) \approx \frac{f(x+3h)-4f(x+2h)+5f(x+h)-5f(x-h)+4f(x-2h)-f(x-3h)}{2h^5},
\qquad O(h^2).
\]

The main conclusion is that the central formulas are more accurate for the same step size because their leading truncation error is $O(h^2)$ rather than $O(h)$.

\subsection{Problem 4(c): error model and optimal step size}

For the forward-difference definition
\[
f'(x) \approx \frac{f(x+h)-f(x)}{h},
\]
the total error is the sum of truncation and roundoff contributions. For $f(x)=\sin x$ at $x=0.5$, the lecture gives the model
\[
g(h)=\left|\frac{h}{2}\sin(0.5)+\frac{\tilde{\varepsilon}}{h}\right|,
\qquad \tilde{\varepsilon}=7\times 10^{-17}.
\]
Balancing the two terms gives
\[
\frac{1}{2}\sin(0.5)-\frac{\tilde{\varepsilon}}{h^2}=0,
\]
so the predicted best step is
\[
h_{\mathrm{best}}
=
\sqrt{\frac{2\tilde{\varepsilon}}{\sin(0.5)}}
\approx 1.71\times 10^{-8}.
\]
This value should be close to the minimum seen in the numerical log-log plot.

\section{Numerical Results}

\subsection{Problem 2: approximating $f'(2)$ for $f(x)=xe^x$}

The exact derivative is
\[
f'(x)=(x+1)e^x,
\qquad
f'(2)=3e^2\approx 22.1671682968.
\]
Using the table values with $h=0.1$, the results are shown in Table~\ref{tab:p2table}.

\begin{table}[H]
\centering
\caption{Problem 2 results from the given table data at $h=0.1$.}
\label{tab:p2table}
\begin{tabular}{lcc}
\toprule
Method & Approximation & Absolute error \\
\midrule
Forward difference & 23.70845000 & 1.54128170 \\
Backward difference & 20.74913000 & 1.41803830 \\
Central difference & 22.22879000 & 0.06162170 \\
Forward 3-point & 22.03231000 & 0.13485830 \\
Backward 3-point & 22.05452500 & 0.11264330 \\
Central 5-point & 22.16699917 & 0.00016913 \\
\bottomrule
\end{tabular}
\end{table}

The central 5-point formula is clearly the best one for the tabulated data. When the step size is swept from $10^{-1}$ to $10^{-12}$, the smallest reported error is $3.55\times 10^{-12}$ for the central 5-point method at $h=10^{-3}$. After that, the error increases again because subtractive cancellation and floating-point roundoff begin to dominate.

\subsection{Problem 3: approximating $f''(2)$ for $f(x)=xe^{x^2}$}

For this function,
\[
f''(x)=(6x+4x^3)e^{x^2},
\qquad
f''(2)=44e^4\approx 2402.3186014583.
\]
The table-data results are listed in Table~\ref{tab:p3table}.

\begin{table}[H]
\centering
\caption{Problem 3 results from the given table data at $h=0.1$.}
\label{tab:p3table}
\begin{tabular}{lcc}
\toprule
Method & Approximation & Absolute error \\
\midrule
Forward 3-point & 4189.712800 & 1787.394199 \\
Backward 3-point & 1468.599900 & 933.718701 \\
Central 3-point & 2460.877300 & 58.558699 \\
Central 5-point & 2399.497458 & 2.821143 \\
\bottomrule
\end{tabular}
\end{table}

Again the central 5-point formula performs best. In the $h$-sweep, its smallest reported error is $6.37\times 10^{-8}$ at $h=10^{-3}$. When $h$ is reduced below about $10^{-4}$, the roundoff term grows rapidly because the formula divides by $h^2$, so very small perturbations in the function values are strongly amplified.

\subsection{Problem 4(a) and 4(b): forward-difference derivative of $f(x)=x(x-1)$}

The analytical derivative is
\[
f'(x)=2x-1,
\qquad
f'(1)=1.
\]
For part (a), the forward-difference program with $\delta=10^{-2}$ gives
\[
f'(1)\approx 1.01,
\]
so the absolute error is exactly $0.01$.

For part (b), the second program evaluates a sequence of $\delta$ values. The results are summarized in Table~\ref{tab:p4b}.

\begin{table}[H]
\centering
\caption{Problem 4(b): derivative accuracy versus $\delta$.}
\label{tab:p4b}
\begin{tabular}{ccc}
\toprule
$\delta$ & Numerical derivative & Absolute error \\
\midrule
$10^{-2}$ & 1.0100000000000009 & $1.000000\times 10^{-2}$ \\
$10^{-4}$ & 1.0000999999998899 & $1.000000\times 10^{-4}$ \\
$10^{-6}$ & 1.0000009999177333 & $9.999177\times 10^{-7}$ \\
$10^{-8}$ & 1.0000000039225287 & $3.922529\times 10^{-9}$ \\
$10^{-10}$ & 1.0000000828403710 & $8.284037\times 10^{-8}$ \\
$10^{-12}$ & 1.0000889005833413 & $8.890058\times 10^{-5}$ \\
$10^{-14}$ & 0.9992007221626509 & $7.992778\times 10^{-4}$ \\
$10^{-16}$ & 0.0000000000000000 & $1.000000\times 10^{0}$ \\
$10^{-18}$ & 0.0000000000000000 & $1.000000\times 10^{0}$ \\
\bottomrule
\end{tabular}
\end{table}

The best value in this list is $\delta=10^{-8}$. The data show the standard pattern very clearly: decreasing $\delta$ first reduces truncation error, but eventually the finite precision of the machine destroys the subtraction $f(x+\delta)-f(x)$.

\subsection{Problem 4(c): reproducing the lecture error plot}

The script \texttt{hw7\_4/c/c.py} reproduces the lecture plot for the forward-difference approximation of $\sin x$ at $x=0.5$. The numerical sweep uses
\[
h=10^{-0.1i},
\qquad i=0,1,\dots,200.
\]
The script reports the best sampled value
\[
h_{\mathrm{best,num}}=10^{-8},
\]
which is consistent with the theoretical estimate
\[
h_{\mathrm{best,theory}}\approx 1.71\times 10^{-8}.
\]
The agreement is good because the sampled grid in $\log_{10}(h)$ has spacing $0.1$, so $10^{-8}$ is the nearest tested point to the theoretical optimum.

\begin{figure}[H]
\centering
\includegraphics[width=0.82\textwidth]{derivative_error_plot_with_theory.png}
\caption{Problem 4(c): numerical forward-difference error and the lecture error model for $f(x)=\sin x$ at $x=0.5$.}
\label{fig:p4c}
\end{figure}

\section{Discussion}

The four tasks illustrate the same practical lesson from different angles. Higher-order symmetric formulas usually outperform one-sided formulas when the same data are available, because their truncation error is smaller. This is why the central 5-point formulas are the best performers in both Problems 2 and 3.

At the same time, the homework also shows why ``smaller step size is always better'' is false. When $h$ or $\delta$ becomes too small, the subtraction of nearly equal floating-point numbers loses significant digits, and division by $h$ or $h^2$ magnifies that loss. The result is an optimal intermediate step size rather than a monotonic improvement.

The selected source files are internally consistent with this interpretation. The electronic derivation in Problem 1, the complete Python workflows in Problems 2 and 3, and the direct lecture-matching scripts in Problem 4 together form a coherent final submission set.

\subsection{Conclusion}

The final conclusions of Homework 7 are:
\begin{enumerate}[label=(\arabic*)]
    \item Central finite-difference formulas are more accurate than forward or backward formulas at the same step size when symmetric data are available.
    \item For Problem 2 and Problem 3, the central 5-point formulas are the most accurate among the tested methods.
    \item Reducing $h$ improves accuracy only up to a point; after that, roundoff error dominates.
    \item In Problem 4(c), the theoretical optimum $h_{\mathrm{best}}\approx 1.71\times 10^{-8}$ agrees well with the numerical minimum near $10^{-8}$.
\end{enumerate}

\vfill
\begin{center}
{\fontsize{28}{32}\selectfont\calligra End of Report}
\end{center}
\vspace*{2em}

\end{document}
